\section{Institutional Background}\label{sec:background}

\subsection{Conditional tokens and binary markets}

A Polymarket binary market resolves to one of two outcomes. The
exchange records the outcomes as ERC-1155 conditional tokens issued
by Gnosis's Conditional Tokens Framework (CTF) on Polygon. Each
market has a YES token and a NO token; the contract guarantees that
one YES + one NO can always be redeemed for one USDC after
resolution. As a consequence, in equilibrium the YES price plus the
NO price equals one and either side can be priced as the
probability of the corresponding outcome.

Liquidity exists for both sides. A trader who is long YES and
believes the implied probability is too high can either sell YES on
the YES book or buy NO on the NO book. The two strategies differ in
inventory effects and slippage but are economically equivalent at
mid. Throughout this paper we report measures on the YES side; YES
and NO books are structural mirrors of each other and the panel
artifacts contain both.

\subsection{CTF Exchange contract architecture}

Trading happens through the Polymarket CTF Exchange smart contract,
which matches signed off-chain orders against resting liquidity and
executes settlement on chain. The contract has two versions. V1
(\texttt{0x4bFb41d5\allowbreak B3570DeFd03\allowbreak C39a9A4D8dE6Bd8\allowbreak B8982E})
ran from launch through the end of April 2026 and is the version our
28-day scrape window targets. V2
(\texttt{0xE111180000d2663C\allowbreak 0091e4f400237545\allowbreak B87B996B})
launched in early April 2026 with a hard cutover; V1 orders were
rejected after the switch. Both versions emit the same
\texttt{OrderFilled} event signature, with topic hash
\texttt{0xd0a08e8c\allowbreak 493f9c94f29311604c9de\allowbreak 1b4e8c8d4c06bd0c789af57f2d65bfec0f6}.

The event payload carries
\texttt{(orderHash},
\texttt{maker}, \texttt{taker},
\texttt{makerAssetId}, \texttt{takerAssetId},
\texttt{makerAmountFilled},
\texttt{takerAmountFilled}, \texttt{fee)}.
The asset ids identify
which side held USDC: \texttt{makerAssetId == 0} means the maker
posted USDC and the taker is the seller of the conditional token;
\texttt{takerAssetId == 0} means the taker posted USDC and is
therefore the buyer. This binary asset-id field is the on-chain
ground truth for trade direction that we use throughout the paper.

\subsection{Polygon block timing and settlement}

Polygon mainnet produces a block roughly every two seconds. Trade
settlement is final to within Heimdall's 256-block reorg buffer,
so we scrape only blocks at depth $\geq 256$ to avoid reorg
contamination. Block timestamps are monotonic but not perfectly
spaced; for our window the empirical mean inter-block interval is
2.0\,s with a small variance, well below the 60-second sample step
used in the panel compute. Where we need to attach a wall-clock
timestamp to an on-chain trade we use linear interpolation from a
single anchor block, which is accurate to within a few seconds over
the 28-day window.

\subsection{Public WebSocket feed structure}

In parallel with on-chain settlement, Polymarket broadcasts a public
WebSocket feed that exposes order-book state to any subscriber. The
feed emits two event types. \texttt{book\_snapshot} carries a
complete L2 snapshot of one side of one market and is sent on
subscription and at irregular intervals thereafter (0.6\% of events
in our archive). \texttt{price\_change} carries a delta against the
last known book state: a single triple
\texttt{(change\_price, change\_side, change\_size)} that updates one
level on one side. A \texttt{change\_size} of zero means the level
is empty; a non-zero size is the new resting quantity at that price.
These events are 99.4\% of the archive.

The feed does not expose taker identity. The \texttt{change\_side}
field labels which side of the book the level update applies to,
not which side initiated the trade that produced the update. This
distinction drives the measurement gap of Section~\ref{sec:limits}:
an inference algorithm working from the feed alone cannot reliably
reconstruct the aggressor sign because the underlying field never
encodes it.

\subsection{Aggressor-sign dependence of standard measures}

Most direction-dependent microstructure measures take the form
\begin{equation*}
  M = \mathbb{E}\!\left[
    \mathrm{sign}_t \cdot
    f(\mathrm{price}_t, \mathrm{mid}_t, \mathrm{size}_t)
  \right],
\end{equation*}
where $\mathrm{sign}_t \in \{-1, +1\}$ is the aggressor sign for
trade $t$ and $f$ is the per-trade contribution. Effective half-spread
($f = \mathrm{price}_t - \mathrm{mid}_t$), realized half-spread,
Kyle's $\lambda$ regression coefficient, and the adverse-selection
component of Glosten-Harris all instantiate this form. A noisy $\mathrm{sign}_t$ does not just attenuate $M$; depending on
the distribution of $f$ across trades, it can flip the sign of $M$
entirely. The 59\% sign-agreement rate of
Section~\ref{ssec:sign-agreement} is what produces the 67\%
effective-spread and 60\% Kyle's $\lambda$ sign-flip rates we report
on the top-100 panel.

\subsection{Comparison to equity-market microstructure}\label{ssec:equity-comparison}

Equity-market benchmarks are what give several of our findings their
scale. \citet{ellis-michaely-ohara-2000} report that the Lee-Ready
algorithm correctly classifies 81\% of trades on Nasdaq, and the
original \citet{lee-ready-1991} study reports a similar accuracy on
NYSE TAQ; the $\sim 59\%$ feed-driven inference rate on Polymarket
is therefore about 22 percentage points below the equity benchmark
on the same algorithm.
Effective half-spreads on liquid US equities post-decimalization sit
in the single-digit-basis-point range
\citep{hasbrouck-2007,foucault-pagano-roell-2013}; on Polymarket,
expressed in basis points relative to mid, the median quoted
half-spread on the central price decile is around 200\,bps, an order
of magnitude wider, consistent with longer-horizon prediction-market
positions and substantially smaller market-maker capital.
Liquidity-provider activity on US equities is heavily concentrated
in a handful of high-frequency firms that account for roughly half
of trading volume \citep{brogaard-hendershott-riordan-2014}; the
median Polymarket market in our panel has $\sim 32$ effective makers
(SF4), suggesting a flatter maker landscape on the venue. Wash
trading on regulated equity exchanges is essentially zero (it is
prohibited under SEC Rule 10b-5); on unregulated cryptocurrency
token exchanges \citet{cong-li-niu-2023} document wash-trade shares
of 25--70\% via a network-classifier approach. Polymarket sits at a
median 1\% under our direct-detection lower bound with a 22\% upper
tail. The token-exchange benchmark is a different incentive
environment (volume-tied listing fees, aggregator-ranking
optimisation, market-making rebates), so we treat 25--70\% as a
sanity bound rather than an apples-to-apples reference. Finally,
collector-side ingestion latency for direct equity feeds is
sub-millisecond and SIP latency is in the hundreds of microseconds
\citep{hasbrouck-2007}; our 41\,ms median for archive ingestion is
two orders of magnitude slower, but that gap reflects our
collector pipeline rather than the exchange.
